\subsection{RAPID routing parameters}
\label{app:routing_parameters}

RAPID routing parameters and parameter-selection rules were shared by the reference network and all simplified variants within each study case. State-specific differences were limited to network geometry, effective routing widths, slopes, and the resulting travel-time parameters. The width-to-depth ratio was estimated once for each case from its reference section and retained across its variants. The routing timestep and spin-up duration were selected separately for each state because both depend on the state-specific distribution of travel-time parameters; they were not independently calibrated.

At bifurcations, discharge was partitioned among downstream routing reaches in proportion to their effective routing widths, with partitioning ratios normalised to sum to one. Mapped reaches were subdivided into RAPID-only subreaches using a target length of 200~m. This subdivision, together with the effective-length floor and stability cap used for $K$ calculation, affected routing resolution and travel-time calculation only; it did not alter the mapped network geometry or geometric analyses.

\begin{table}[t]
\centering
\caption{RAPID routing parameters and parameter-selection rules. Values labelled as ranges summarise the currently completed anchored runs and should be verified after all final runs have been completed.}
\label{tab:routing_parameters}
\begin{tabular}{p{0.31\linewidth}p{0.13\linewidth}p{0.20\linewidth}p{0.26\linewidth}}
\hline
Parameter & Symbol & Value & Description \\
\hline
Muskingum weighting factor & $X$ & 0.1 & Applied to all cases and states. \\
Manning roughness & $n$ & 0.35 & Applied to all cases and states. \\
Width-to-depth ratio & $k_b$ & 120.1--314.5 & Case-specific value estimated from the reference section using BASED; held constant across variants within a case. \\
Minimum celerity & $c_{\min}$ & 0.05--0.07~m~s$^{-1}$ & Case-specific lower celerity cap selected for numerical stability. \\
Maximum celerity & $c_{\max}$ & 2.524~m~s$^{-1}$ & Upper celerity cap applied to all cases and states. \\
Target subreach length & $T_l$ & 200~m & Target length used to create RAPID-only subreaches. \\
Effective-length floor & $L_{\mathrm{eff,min}}$ & 200~m & Minimum length used in travel-time parameter calculation. \\
Stability effective-length cap & --- & Enabled & Applied near the lower celerity bound; 10~s margin and 0.03 relative celerity buffer. \\
Routing timestep & $\Delta t$ & 620--4,200~s & Selected per state from the allowable Muskingum stability interval. \\
Spin-up multiplier & $m_{\mathrm{spinup}}$ & 1.1 & Spin-up duration equals 1.1 times the longest downstream $K$ path for that state; an equally long terminal period was added. \\
\hline
\end{tabular}
\end{table}

% FINAL-RUN CHECK: Recalculate the reported ranges for k_b, c_min, and Delta t from all completed final anchored experiments, then remove the qualification in the table caption.
